\section{Literature Survey and Related Work}
\label{ch:literature}

\subsection{Overview}

Work relevant to this project falls into three strands: diagnostics built
directly on the OBD-II interface, machine learning for automotive predictive
maintenance in general, and the specific model family (gradient-boosted
trees with attribution methods) that we deploy. Reading across the three
strands, a consistent pattern emerges: the studies with realistic fault data
do not use the consumer OBD-II interface, and the studies on the consumer
interface do not have realistic fault data. Our project sits precisely in
that intersection, which is why the survey below pays as much attention to
each paper's \emph{data source} as to its model.

\subsection{Related Work}

\subsubsection{Diagnostics on the OBD-II interface}

Rimpas et al.~\cite{rimpas2020} logged OBD-II parameters --- fuel trims,
coolant temperature, ignition timing, and airflow among them --- to monitor
vehicle operation and fuel consumption, and discussed how deviations in
individual PIDs indicate specific degradations. Their work confirms the
diagnostic value of exactly the PIDs our system builds features on, but the
analysis is per-parameter and threshold-style. No classifier is trained, and
no labeled faults are used; the interpretation of a drifting trim value is
left to the reading engineer.

Shafi et al.~\cite{shafi2018} built a remote health-monitoring system with
an OBD-II dongle streaming to a phone and applied several classifiers
(decision trees, SVM, k-NN) to detect anomalies in four vehicle subsystems.
Their fault examples, however, come from threshold rules applied to the same
stream the classifiers read --- the labels are derived from the data rather
than from independently verified physical faults. The system demonstrates
the plumbing of phone-based OBD-II ML but not validated fault detection.

The carOBD dataset and its accompanying thesis by
Maranh\~ao~\cite{maranhao_carobd} provide the raw material our project
trains on: 129 sessions of 1\,Hz OBD-II logging from a single Toyota Etios,
covering idle, commute, campus, highway, and long-trip driving. The dataset
contains healthy driving only. To our knowledge no public OBD-II corpus at
this rate contains verified, labeled faults; we searched for one and did not
find it, and this absence is the central constraint the project designs
around.

That absence is structural, not an oversight anyone will quickly correct.
Collecting labeled fault data at scale means deliberately breaking cars, or
waiting for a fleet's worth of them to break naturally while logging
continuously, and then joining the logs to verified repair records. The first is expensive and, done carelessly,
dangerous; the second requires fleet access and years of patience. Whoever
pays for either has a commercial reason to keep the result private. The
public-data vacuum that Theissler et al.\ describe at the field level is, in
other words, structural, and a graduation project cannot fill it by trying
harder. It can only decide how to work honestly inside it, which is the
decision Section~\ref{sec:injector} documents.

\subsubsection{Machine learning for automotive predictive maintenance}

Theissler et al.~\cite{theissler2021} survey machine learning for predictive
maintenance across the automotive industry and catalogue the field's
recurring obstacles: scarce labeled failure data, class imbalance, and the
gap between lab results and fleet deployment. They note explicitly that
run-to-failure data for vehicles is rarely public because it is expensive
and commercially sensitive --- the same wall we hit, met in their survey as
a field-wide condition rather than a local complaint.

Prytz et al.~\cite{prytz2015} predicted compressor repairs in a commercial
truck fleet from logged vehicle data and workshop maintenance records. Their
labels are real --- actual repairs --- but the data source is an OEM fleet
telematics database with service-record joins, not the standardized consumer
port, and the features are aggregated per vehicle-visit rather than
streaming windows. The study demonstrates that drift-before-failure is
learnable in principle, at fleet scale, with access no individual owner has.

Theissler~\cite{theissler2017} addressed the known/unknown fault problem
with ensemble anomaly detection over in-vehicle recordings, arguing that a
detector must handle fault types absent from training. Our two-pillar
design --- a supervised classifier for the four modeled faults plus a
one-class IsolationForest~\cite{liu2008} trained on healthy windows only ---
is a direct application of that argument at consumer-hardware scale.

Vergara et al.~\cite{vergara2023} published EngineFaultDB, a labeled engine
fault dataset with baseline classification results. The faults are real, but
they were produced on a laboratory engine test bench with instrumentation
beyond the OBD-II PID set, sampled under controlled load profiles. We
considered benchmarking against it and declined: a bench dynamometer corpus
and a 1\,Hz road-driving corpus differ in sensors, dynamics, and noise floor
to the point where a comparison table would mislead more than inform.

\subsubsection{Models and attribution}

Gradient-boosted decision trees remain the strongest general-purpose family
for tabular sensor features, and XGBoost~\cite{chen2016} is its standard
implementation; random forests~\cite{breiman2001} provide the natural
bagging baseline. For attribution we use SHAP~\cite{lundberg2017}, whose
TreeExplainer computes exact Shapley values for tree ensembles --- on a
workshop floor, ``LTFT pushed this prediction toward fuel-system by 0.43''
is an explanation a mechanic can act on. The protocol layer underneath the
whole pipeline is standardized by SAE J1979~\cite{saej1979} and
ISO 15765-4~\cite{iso15765}, and our tooling builds on the python-obd
library~\cite{pythonobd}, Streamlit~\cite{streamlit}, and
FastAPI~\cite{fastapi}.

Table~\ref{tab:relwork} lays the surveyed work along the two axes that
decide whether a result transfers to our setting: the realism of its fault
labels and the interface its data came through. The pattern is stark ---
every row has one or the other, never both.

\begin{table}[H]
  \centering
  \caption{Related work along the two axes that matter for our problem:
    fault-label realism and data interface. No surveyed work combines
    verified faults with the 1\,Hz consumer OBD-II interface.}
  \label{tab:relwork}
  \begin{tabular}{lll}
    \toprule
    Work & Fault labels & Data source\\
    \midrule
    Rimpas et al.~\cite{rimpas2020} & None (per-PID thresholds) & Consumer OBD-II\\
    Shafi et al.~\cite{shafi2018} & Derived from the stream & Consumer OBD-II\\
    Prytz et al.~\cite{prytz2015} & Real (repair records) & OEM fleet telematics\\
    Vergara et al.~\cite{vergara2023} & Real (induced) & Bench dynamometer\\
    Theissler~\cite{theissler2017} & Mixed known/unknown & In-vehicle recordings\\
    \textbf{This work} & Synthetic, physics-constrained & Consumer OBD-II, 1\,Hz\\
    \bottomrule
  \end{tabular}
\end{table}

\subsection{Gap Analysis}

Laying the surveyed work side by side exposes four gaps.

\textbf{Realistic faults and consumer hardware do not co-occur.} Bench
studies such as EngineFaultDB~\cite{vergara2023} have verified faults but
non-OBD instrumentation and dynamometer conditions. Fleet studies such as
Prytz et al.~\cite{prytz2015} have real repair labels but proprietary
telematics access. OBD-II studies such as Shafi et al.~\cite{shafi2018}
run on the consumer port but lack independently verified fault labels.
No surveyed work trains and evaluates on the 1\,Hz consumer interface with
trustworthy fault ground truth --- because the data to do so does not
publicly exist.

\textbf{Single-vehicle scope is the norm, and transfer is untested.} The
carOBD corpus is one Etios; EngineFaultDB is one bench engine; most OBD-II
papers instrument one car. Almost none of the surveyed systems state what
happens when the trained model meets a different vehicle --- different
idle speed, different trim baseline, different PID set. In our experience
that encounter is where deployments actually fail.

\textbf{Missing sensors are not handled honestly.} Every surveyed system
assumes its input channels exist. But PID support varies car to car: our
own test vehicle exposes no MAP sensor, the primary channel for one of our
four fault classes. No prior system we found defines what it reports when
the evidence for a fault class is physically absent.

\textbf{Synthetic injection, where used, is not physics-audited.} Adding
offsets or noise to sensor channels is common in anomaly detection
research, but the surveyed automotive works do not constrain injected
values to ECU causality --- fuel trims that lead instead of react, coolant
that jumps faster than thermal mass allows, MAP exceeding barometric
pressure. Models trained on physically impossible signatures learn
artifacts that no real fault will ever reproduce.

\subsection{How This Project Fills the Gap}

Against the first gap, we accept the absence of real-fault data as a design
input rather than pretending otherwise: we train on physics-constrained
synthetic faults injected into real 1\,Hz driving data, we measure how much
the scores depend on the injector's own coefficients (the withheld-parameter
experiment of Section~\ref{sec:evalmethod}), and we label every reported
number as synthetic self-consistency. The deliverable is a detection
algorithm plus a collection protocol, and the report never claims more.

Against the second, every feature is z-scored against the source vehicle's
own healthy baseline (Section~\ref{sec:features}), so the model reasons in
``standard deviations from this car's normal'' rather than raw units, and
the deployed pipeline re-fits that baseline per vehicle. Our evaluation
splits by session and by session \emph{family}, the closest available proxy
for context transfer within a single-vehicle corpus.

Against the third, the pipeline carries an explicit evaluability contract:
each fault class declares the PIDs its primary signature requires, and a
vehicle that lacks them gets ``Untested'' for that class instead of a score
(Section~\ref{sec:dashboard}). Chapter~\ref{ch:results} shows this contract
working on a real MAF-based car.

Against the fourth, the injection engine enforces hard physical clamps and
documented ECU response dynamics --- including the speed-density vacuum-leak
behavior and the STFT-to-LTFT handoff --- with every injection reproducible
from stored parameters (Section~\ref{sec:injector}). The physics arguments
are written into the source alongside the code they constrain.
